\documentclass[12pt,a4paper]{article}

\usepackage[T2A]{fontenc}
\usepackage[utf8]{inputenc}
\usepackage[russian]{babel}
\usepackage{amsmath,amssymb,amsthm,mathtools}
\usepackage{geometry}
\usepackage{hyperref}
\usepackage{enumitem}
\usepackage{array}
\usepackage{longtable}
\usepackage{stmaryrd}
\geometry{margin=2.2cm}

\title{Расширенный ответ по теории алгоритмов\\
для вступительных экзаменов в аспирантуру}
\author{}
\date{}

\newtheorem{definition}{Определение}
\newtheorem{theorem}{Теорема}
\newtheorem{proposition}{Утверждение}
\newtheorem{corollary}{Следствие}
\newtheorem{example}{Пример}
\newtheorem{remark}{Замечание}
\newtheorem{lemma}{Лемма}


\title{Билет №3 \\ \small Понятие алгоритма. Разрешимые и перечислимые языки. Алгоритмически неразрешимые задачи. Проблема останова. Теорема Райса. (ИТМО) \\ Понятие алгоритма и его уточнения: машины Тьюринга, нормальные алгоритмы Маркова, рекурсивные функции. Эквивалентность данных формальных моделей алгоритмов. Понятие об алгоритмической неразрешимости. Примеры алгоритмически неразрешимых проблем. (СПбГУ) \\ Понятие алгоритма и его уточнения: машины Тьюринга, нормальные алгоритмы Маркова, рекурсивные функции. Понятие об алгоритмической неразрешимости. Примеры алгоритмически неразрешимых проблем (СПбПУ)}
\author{Сериков Владислав Александрович}
\date{23 апреля 2026}

\begin{document}

\maketitle
\tableofcontents
\newpage

\section{Понятие алгоритма и вычисления}

\subsection{Интуитивное понятие алгоритма}
Центральным понятием курса является не столько \emph{алгоритм}, сколько \emph{вычисление}. Алгоритм --- это способ описания вычисления, то есть точное указание, \emph{как} нужно действовать, чтобы получить результат.

\begin{definition}
Алгоритм --- это точное предписание, задающее конечный или потенциально бесконечный процесс последовательного выполнения элементарных шагов.
\end{definition}

\subsection{Почему понадобилось строгое определение}
Исторически формализация понятия алгоритма понадобилась потому, что математики захотели отвечать не только на вопрос \emph{``как решить задачу?''}, но и на вопрос \emph{``а вообще можно ли её решить алгоритмом?''}.

Ключевую роль здесь сыграли:
\begin{itemize}
    \item программа Гильберта;
    \item проблема разрешения (\emph{Entscheidungsproblem});
    \item более общий интерес к существованию механических процедур доказательства и вычисления.
\end{itemize}

Чтобы доказать, что некоторого алгоритма \emph{не существует}, надо сначала чётко определить, что считается алгоритмом. Именно отсюда возникают \textbf{формальные модели алгоритмов}.

\subsection{Вычисление как функция и как распознавание языка}
Выделяются две вычислительные парадигмы:
\begin{enumerate}
    \item вычисление как \textbf{числовая функция};
    \item вычисление как \textbf{преобразование и распознавание слов}.
\end{enumerate}

Обе парадигмы эквивалентны по выразительной силе, так как числа можно кодировать словами, а слова --- числами.

\begin{definition}
Алфавит $\Sigma$ --- это конечное множество символов.
\end{definition}

\begin{definition}
Слово в алфавите $\Sigma$ --- конечная последовательность символов из $\Sigma$.
Множество всех слов обозначается $\Sigma^*$.
\end{definition}

\begin{definition}
Язык в алфавите $\Sigma$ --- произвольное подмножество $\Sigma^*$.
\end{definition}

\begin{definition}
Задача распознавания языка $L \subseteq \Sigma^*$ состоит в следующем: по слову $w\in\Sigma^*$ определить, принадлежит ли оно языку $L$.
\end{definition}

\begin{example}
Пусть $\Sigma=\{| \}$ и
\[
P=\{\,|^n \mid n \text{ --- простое число}\,\}.
\]
Тогда задача распознавания языка $P$ эквивалентна задаче проверки простоты натурального числа.
\end{example}

\section{Формальные модели алгоритмов}
Существует несколько эквивалентных моделей вычисления, например:
\begin{itemize}
    \item нормальные алгоритмы Маркова;
    \item машины Тьюринга;
    \item рекурсивные функции.
\end{itemize}

Эквивалентность этих моделей позволяет говорить о едином математическом понятии вычислимости.

\section{Нормальные алгоритмы Маркова}

\subsection{Исходные определения}

\begin{definition}
Пусть $\Sigma$ --- алфавит. \textbf{Марковской подстановкой} называется пара слов
\[
P \to Q.
\]
\end{definition}

\begin{definition}
Подстановка $P\to Q$ \textbf{применима} к слову $M$, если $P$ является подсловом $M$.
\end{definition}

\begin{definition}
\textbf{Применение подстановки} к слову $M$ --- это замена \emph{первого} вхождения $P$ в $M$ на $Q$.
\end{definition}

\begin{definition}
\textbf{Заключительная подстановка} записывается в виде
\[
P\to. Q
\]
и означает, что после её применения работа алгоритма должна завершиться.
\end{definition}

\begin{definition}
\textbf{Схема нормального алгоритма} --- это конечный упорядоченный список подстановок
\[
P_1\to(.)Q_1,\quad P_2\to(.)Q_2,\quad \dots,\quad P_n\to(.)Q_n.
\]
\end{definition}

\begin{definition}
\textbf{Однократное применение схемы} к слову --- это применение первой по порядку подстановки, которая применима к данному слову.
\end{definition}

\begin{definition}
\textbf{Нормальный алгоритм} --- это последовательность однократных применений схемы к начальному слову, продолжающаяся до тех пор, пока:
\begin{enumerate}
    \item либо не будет применена заключительная подстановка;
    \item либо ни одна подстановка не окажется применимой.
\end{enumerate}
\end{definition}

\subsection{Алгоритм работы: шаги}
Для любого нормального алгоритма шаги вычисления можно выписать буквально:

\begin{enumerate}
    \item Взять текущее слово $W$.
    \item Просмотреть подстановки схемы сверху вниз.
    \item Найти первую подстановку, применимую к $W$.
    \item Заменить \emph{первое} вхождение её левой части на правую часть.
    \item Если подстановка заключительная, остановиться.
    \item Если применимых подстановок нет, остановиться.
    \item Иначе перейти к следующему шагу.
\end{enumerate}

\subsection{Иллюстрация на простом примере}
Рассмотрим схему
\[
\left\{
\begin{array}{l}
ab\to \epsilon,\\
a\to b.
\end{array}
\right.
\]

Применим её к слову $aababab$:
\[
aababab \Rightarrow aabab \Rightarrow aab \Rightarrow a \Rightarrow b.
\]

Разбор по шагам:
\begin{enumerate}
    \item В слове $aababab$ есть первое вхождение $ab$; применяем $ab\to \epsilon$:
    \[
    aababab \Rightarrow aabab.
    \]
    \item Аналогично:
    \[
    aabab \Rightarrow aab.
    \]
    \item Снова:
    \[
    aab \Rightarrow a.
    \]
    \item Теперь $ab$ уже нет, но есть $a$, поэтому применяем вторую подстановку:
    \[
    a \Rightarrow b.
    \]
    \item Для слова $b$ ни одна подстановка не применима, процесс завершается.
\end{enumerate}

\subsection{Пример вычисления функции: увеличение унарного числа}
Пусть натуральное число $n$ кодируется словом $|^n$. Рассмотрим схему:
\[
\{| \to. ||\}.
\]

\paragraph{Смысл.} Она увеличивает число на единицу.

\paragraph{Шаги алгоритма:}
\begin{enumerate}
    \item Найти первый символ $|$.
    \item Заменить его на $||$.
    \item Так как подстановка заключительная, завершить вычисление.
\end{enumerate}

\paragraph{Пример:}
\[
|| \Rightarrow |||.
\]

\paragraph{Доказательство корректности.}
Пусть входное слово имеет вид $|^n$, $n\ge 1$. Тогда левая часть подстановки присутствует в слове. После применения подстановки общее число символов увеличивается на $1$, так как один символ заменяется двумя. Следовательно, результат имеет вид $|^{n+1}$. Это и означает прибавление единицы.

\subsection{Пример: вычитание единицы}
Схема:
\[
\{| \to. \varepsilon\}.
\]

\paragraph{Шаги:}
\begin{enumerate}
    \item Найти первый символ $|$.
    \item Удалить его.
    \item Остановиться.
\end{enumerate}

\paragraph{Пример:}
\[
||| \Rightarrow ||.
\]

\paragraph{Корректность.}
Если вход есть $|^n$ при $n\ge 1$, то результатом становится $|^{n-1}$. Для пустого слова схема неприменима; в таком виде функция определена только на непустых словах.

\subsection{Пример проверки делимости на три}
В конспекте приведена схема:
\[
\left\{
\begin{array}{l}
||| \to \varepsilon,\\
|| \to. \varepsilon,\\
| \to. \varepsilon,\\
\varepsilon \to. |
\end{array}
\right.
\]

\paragraph{Идея.}
Пока можно, удаляются тройки палочек. После этого остаётся слово длины $0$, $1$ или $2$.
\begin{itemize}
    \item Если осталось $0$, то применяется $\varepsilon\to. |$, и результатом становится $|$.
    \item Если осталось $1$ или $2$, они стираются заключительной подстановкой, и результатом становится $\varepsilon$.
\end{itemize}

\paragraph{Пример 1: $7$ не делится на $3$}
\[
||||||| \Rightarrow |||| \Rightarrow | \Rightarrow \varepsilon.
\]

\paragraph{Пример 2: $9$ делится на $3$}
\[
||||||||| \Rightarrow |||||| \Rightarrow ||| \Rightarrow \varepsilon \Rightarrow |.
\]

\paragraph{Доказательство корректности.}
Пусть входное слово есть $|^n$.
После конечного числа применений первой подстановки длина уменьшается на $3$ на каждом шаге, поэтому алгоритм приходит к остатку длины $r$, где $r\in\{0,1,2\}$ и
\[
n\equiv r \pmod 3.
\]
Дальше:
\begin{itemize}
    \item если $r=0$, то применима только подстановка $\varepsilon\to. |$, и результат равен $|$;
    \item если $r=1$ или $r=2$, то применима одна из заключительных подстановок $| \to. \varepsilon$ или $||\to.\varepsilon$, и результат равен $\varepsilon$.
\end{itemize}
Следовательно, алгоритм корректно различает числа, делящиеся и не делящиеся на $3$.

\subsection{Алгоритмы над расширенным алфавитом}
Иногда для решения задачи вводят \emph{вспомогательные символы}. Это называется работой алгоритма \textbf{над} алфавитом.

\paragraph{Пример: удалить первую букву слова.}
Для алфавита $\Sigma=\{a,b\}$ предлагается расширить алфавит символом $*$:
\[
\left\{
\begin{array}{l}
*a \to. \varepsilon,\\
*b \to. \varepsilon,\\
\varepsilon \to *
\end{array}
\right.
\]

\paragraph{Шаги алгоритма:}
\begin{enumerate}
    \item В начало слова вставляется служебный символ $*$.
    \item Если после него стоит $a$, применяется $*a\to.\varepsilon$.
    \item Если после него стоит $b$, применяется $*b\to.\varepsilon$.
\end{enumerate}

\paragraph{Пример:}
\[
abab \Rightarrow *abab \Rightarrow bab.
\]

\paragraph{Корректность.}
Символ $*$ маркирует начало слова. После его вставки первая буква оказывается прямо справа от него. Заключительная подстановка удаляет сразу пару ``служебный символ + первая буква'', после чего работа завершается. Следовательно, удаляется ровно первый символ слова.

\subsection{Нормально вычислимые функции}
\begin{definition}
Функция $f$, заданная на словах алфавита $\Sigma$, называется \textbf{нормально вычислимой}, если существует нормальный алгоритм над $\Sigma$, который каждое слово $P$ из области определения переводит в слово $f(P)$.
\end{definition}

Если на некотором входе алгоритм не завершается, считается, что функция на этом входе не определена.

\section{Машины Тьюринга}

\subsection{Интуитивное описание}
Машина Тьюринга --- это абстрактное вычислительное устройство, состоящее из:
\begin{itemize}
    \item бесконечной ленты;
    \item считывающей/записывающей головки;
    \item устройства управления с конечным числом состояний.
\end{itemize}

Она предназначена для максимально простой, но универсальной формализации вычисления.

\subsection{Формальное определение}
\begin{definition}
Машина Тьюринга задаётся семёркой
\[
M=(Q,\Sigma,\Gamma,\delta,q_0,B,F),
\]
где:
\begin{itemize}
    \item $Q$ --- конечное множество состояний;
    \item $\Sigma$ --- входной алфавит;
    \item $\Gamma$ --- ленточный алфавит, причём $\Sigma\subseteq\Gamma$;
    \item $q_0\in Q$ --- начальное состояние;
    \item $B\in \Gamma$ --- пробельный символ;
    \item $F\subseteq Q$ --- множество допускающих состояний;
    \item $\delta:Q\times \Gamma \to Q\times \Gamma \times \{L,R\}$ --- функция переходов.
\end{itemize}
\end{definition}

\subsection{Шаги работы машины Тьюринга}
На одном шаге машина:
\begin{enumerate}
    \item считывает символ в текущей клетке;
    \item смотрит текущее состояние;
    \item по функции $\delta$ определяет:
    \begin{itemize}
        \item новое состояние;
        \item символ, который нужно записать;
        \item направление движения $L$ или $R$;
    \end{itemize}
    \item записывает новый символ;
    \item сдвигает головку;
    \item переходит в новое состояние.
\end{enumerate}

\subsection{Остановка}
Машина останавливается, если:
\begin{itemize}
    \item либо попадает в допускающее состояние;
    \item либо функция переходов не определена для текущей пары ``состояние--символ''.
\end{itemize}

\subsection{Конфигурация}
\begin{definition}
\textbf{Конфигурация} машины Тьюринга --- это описание текущего состояния вычисления: содержимого ленты, положения головки и текущего состояния устройства управления.
\end{definition}

\subsection{Минимальный пример}
Рассмотрим машину:
\[
M=(\{q_0,q_1\},\{0,1\},\{0,1,B\},\delta,q_0,B,\{q_1\}),
\]
где
\[
\delta(q_0,1)=(q_0,1,R),\qquad
\delta(q_0,0)=(q_1,1,R).
\]

\paragraph{Смысл.}
Машина идёт вправо по символам $1$. Как только встретит первый $0$, она заменит его на $1$, перейдёт в допускающее состояние $q_1$ и остановится.

\paragraph{Пошаговый пример на слове $111011$:}
\begin{enumerate}
    \item в состоянии $q_0$ читает $1$, остаётся в $q_0$, идёт вправо;
    \item снова читает $1$, остаётся в $q_0$, идёт вправо;
    \item снова читает $1$, остаётся в $q_0$, идёт вправо;
    \item читает $0$, записывает $1$, переходит в $q_1$, идёт вправо;
    \item машина в недопустимом состоянии, работа завершена.
\end{enumerate}

\paragraph{Результат преобразования:}
\[
111011 \mapsto 111111.
\]

\subsection{Распознавание языка машиной Тьюринга}
\begin{definition}
Языком машины Тьюринга называется множество слов, на которых машина за конечное число шагов попадает в допускающее состояние:
\[
L(M)=\{\,w\in\Sigma^* \mid q_0w \vdash_M^* \alpha p\beta,\ p\in F\,\}.
\]
\end{definition}

\paragraph{Язык предыдущей машины.}
\[
L(M)=\{\,w\in\{0,1\}^* \mid \text{в }w\text{ есть хотя бы один }0\,\}.
\]

\paragraph{Доказательство корректности.}
\begin{itemize}
    \item Если во входном слове есть хотя бы один ноль, то при движении вправо машина обязательно когда-нибудь впервые встретит ноль и перейдёт в состояние $q_1$, то есть допустит слово.
    \item Если же нуля нет, то машина пройдёт по всем единицам, дойдёт до пробела $B$, а переход для $(q_0,B)$ не определён; она остановится не в допускающем состоянии.
\end{itemize}
Следовательно, распознаваемый язык описан верно.

\subsection{Машина как преобразователь}
Если машина используется не для ответа ``да/нет'', а для получения выходного слова, то важен результат на ленте после окончания работы. Тогда она рассматривается как \textbf{машина-преобразователь}.

\section{Рекурсивные функции}

\subsection{Общая идея}
Третья модель алгоритма --- рекурсивные функции. Они строятся через:
\begin{itemize}
    \item инициальные функции;
    \item композицию;
    \item примитивную рекурсию;
    \item неограниченный поиск.
\end{itemize}

\subsection{Частично вычислимые и вычислимые функции}
\begin{definition}
Функция называется \textbf{частично вычислимой} (частично рекурсивной), если существует регистровая машина, вычисляющая её.
\end{definition}

\begin{definition}
Частично вычислимая функция, определённая на всех входах, называется \textbf{вычислимой} (рекурсивной).
\end{definition}

Обозначения:
\[
\mathcal P \quad \text{--- множество частично рекурсивных функций,}
\]
\[
\mathcal R \quad \text{--- множество рекурсивных функций.}
\]

\subsection{Инициальные функции}
\begin{definition}
Инициальные функции:
\begin{enumerate}
    \item нуль-функция
    \[
    Z(x)=0;
    \]
    \item функция следования
    \[
    S(x)=x+1;
    \]
    \item функции-проекторы
    \[
    P_i^n(x_1,\dots,x_n)=x_i.
    \]
\end{enumerate}
\end{definition}

\subsection{Композиция}
\begin{definition}
Если заданы функции
\[
f(x_1,\dots,x_n), \qquad g_1(z),\dots,g_n(z),
\]
то их композиция:
\[
h(z)=f(g_1(z),\dots,g_n(z)).
\]
\end{definition}

\paragraph{Шаги алгоритма композиции:}
\begin{enumerate}
    \item по входу $z$ вычислить все промежуточные значения $g_1(z),\dots,g_n(z)$;
    \item подать их на вход функции $f$;
    \item вычислить результат $f(g_1(z),\dots,g_n(z))$.
\end{enumerate}

\paragraph{Минимальный пример.}
\[
f(x,y)=x+y,\qquad g(z)=z+1.
\]
Тогда
\[
h(z)=f(g(z),g(z))=(z+1)+(z+1)=2z+2.
\]

\paragraph{Корректность.}
Если все функции $f,g_1,\dots,g_n$ вычислимы, то их вычисления заканчиваются. Поэтому можно последовательно получить все промежуточные значения и вычислить итог. Следовательно, композиция сохраняет вычислимость.

\subsection{Примитивная рекурсия}
\begin{definition}
Функция $f(x,y)$ определяется \textbf{примитивной рекурсией} через функции $h(y)$ и $g(x,y,z)$, если
\[
f(0,y)=h(y),
\]
\[
f(x+1,y)=g(x,y,f(x,y)).
\]
\end{definition}

\paragraph{Шаги алгоритма примитивной рекурсии:}
\begin{enumerate}
    \item если первый аргумент равен $0$, вернуть $h(y)$;
    \item иначе вычислить предыдущее значение $f(x,y)$;
    \item применить функцию $g$ к тройке $(x,y,f(x,y))$;
    \item полученное значение объявить результатом.
\end{enumerate}

\subsection{Пример: сложение}
Определим сложение:
\[
add(0,y)=y,
\]
\[
add(x+1,y)=S(add(x,y)).
\]

\paragraph{Вычисление $add(3,2)$:}
\[
add(0,2)=2,
\]
\[
add(1,2)=S(add(0,2))=3,
\]
\[
add(2,2)=S(add(1,2))=4,
\]
\[
add(3,2)=S(add(2,2))=5.
\]

\paragraph{Доказательство корректности.}
Докажем по индукции по $x$, что
\[
add(x,y)=x+y.
\]

\textbf{База:}
\[
add(0,y)=y=0+y.
\]

\textbf{Переход:}
предположим, что
\[
add(x,y)=x+y.
\]
Тогда
\[
add(x+1,y)=S(add(x,y))=add(x,y)+1=(x+y)+1=(x+1)+y.
\]
Следовательно, формула верна для всех $x$.

\subsection{Пример: умножение}
Определим умножение:
\[
mult(0,y)=0,
\]
\[
mult(x+1,y)=add(mult(x,y),y).
\]

\paragraph{Пример:}
\[
mult(2,3)=6.
\]

\paragraph{Идея корректности.}
Это просто повторное прибавление числа $y$ к самому себе $x$ раз. Индукция по $x$ полностью аналогична предыдущему случаю.

\subsection{Неограниченный поиск}
\begin{definition}
Если задана функция $g(x,y)$, то \textbf{оператор неограниченного поиска} определяет
\[
f(y)=\mu x\,[g(x,y)=0],
\]
то есть наименьшее $x$, для которого $g(x,y)=0$.
\end{definition}

\paragraph{Шаги алгоритма:}
\begin{enumerate}
    \item положить $x:=0$;
    \item вычислить $g(x,y)$;
    \item если результат равен $0$, вернуть $x$;
    \item иначе увеличить $x$ на $1$ и повторить.
\end{enumerate}

\paragraph{Минимальный пример.}
Пусть
\[
g(x,y)=|x-y|.
\]
Тогда
\[
\mu x\,[g(x,y)=0]=y.
\]

\paragraph{Замечание.}
Этот алгоритм может не остановиться, если подходящего $x$ не существует. Поэтому здесь естественно возникают \emph{частично} рекурсивные функции.

\subsection{Примитивно рекурсивные функции}
\begin{definition}
Множество примитивно рекурсивных функций --- это наименьшее множество, содержащее инициальные функции и замкнутое относительно композиции и примитивной рекурсии.
\end{definition}

\subsection{Частично рекурсивные функции}
\begin{theorem}
Множество частично рекурсивных функций $\mathcal P$ есть замыкание множества инициальных функций относительно:
\begin{itemize}
    \item композиции;
    \item примитивной рекурсии;
    \item неограниченного поиска.
\end{itemize}
\end{theorem}

\section{Эквивалентность моделей алгоритмов}

\subsection{Формулировка}
\begin{theorem}
Нормальные алгоритмы Маркова, машины Тьюринга и рекурсивные функции эквивалентны по вычислительной силе.
\end{theorem}

\subsection{Что означает эквивалентность}
Эквивалентность означает:
\begin{quote}
любая функция или задача, вычислимая в одной модели, вычислима и в любой другой.
\end{quote}

То есть:
\begin{itemize}
    \item если функция вычислима машиной Тьюринга, то существует эквивалентный нормальный алгоритм Маркова;
    \item если функция рекурсивна, то её можно реализовать машиной Тьюринга;
    \item если функция нормально вычислима, то она частично рекурсивна;
    \item и так далее.
\end{itemize}

\subsection{Почему это принципиально важно}
Именно эквивалентность моделей оправдывает обсуждение \textbf{неразрешимости}:
если задача не решается в одной модели, то она не решается ни в какой другой разумной формальной модели алгоритма.

\subsection{Тезис Чёрча--Тьюринга}
\begin{definition}
\textbf{Тезис Чёрча--Тьюринга} утверждает: всякая интуитивно вычислимая функция вычислима машиной Тьюринга.
\end{definition}

\begin{remark}
Это не теорема в строгом математическом смысле, а фундаментальный тезис, связывающий формальное понятие вычислимости с интуитивным.
\end{remark}

\section{Разрешимые и перечислимые языки}

\subsection{Разрешимые языки}
\begin{definition}
Язык $L\subseteq\Sigma^*$ называется \textbf{разрешимым} (рекурсивным), если существует алгоритм, который для любого входного слова $w$ за конечное время отвечает, принадлежит ли $w$ языку $L$.
\end{definition}

Эквивалентно: существует машина Тьюринга, которая на любом входе останавливается и допускает в точности слова языка $L$.

\subsection{Перечислимые языки}
\begin{definition}
Язык $L$ называется \textbf{перечислимым} (рекурсивно перечислимым, полуразрешимым), если существует алгоритм, который:
\begin{itemize}
    \item для слов $w\in L$ останавливается и подтверждает принадлежность;
    \item для слов $w\notin L$ может не останавливаться.
\end{itemize}
\end{definition}

\subsection{Эквивалентное понимание перечислимости}
Язык рекурсивно перечислим тогда и только тогда, когда существует процедура, которая со временем может перечислить все его слова, возможно с повторениями.

\subsection{Соотношения между классами языков}
\begin{theorem}
Каждый разрешимый язык перечислим.
\end{theorem}

\begin{proof}
Если есть алгоритм, всегда завершающийся с ответом ``да/нет'', то его можно использовать как полурешатель. На словах языка он останавливается и подтверждает принадлежность.
\end{proof}

\begin{theorem}
Если язык $L$ и его дополнение $\overline L$ перечислимы, то язык $L$ разрешим.
\end{theorem}

\begin{proof}
Запустим одновременно два полурешателя:
\begin{itemize}
    \item один для $L$;
    \item другой для $\overline L$.
\end{itemize}
Для любого слова $w$ ровно одно из утверждений истинно:
\[
w\in L \quad \text{или} \quad w\in \overline L.
\]
Следовательно, один из двух процессов когда-нибудь завершится. Это и даёт полный решатель.
\end{proof}

\subsection{Содержательная интерпретация}
\begin{itemize}
    \item \textbf{разрешимость} означает наличие полного алгоритма;
    \item \textbf{перечислимость} означает возможность эффективно распознать ``да'', но не обязательно ``нет'';
    \item \textbf{неперечислимость} означает отсутствие даже такого одностороннего распознавания.
\end{itemize}

\section{Полурекурсивные отношения и полуразрешимость}

\subsection{Определение}
\begin{definition}
Отношение $P(x)$ называется \textbf{полурекурсивным}, если существует частично рекурсивная функция $f$ такая, что
\[
P(x)\iff f(x)\downarrow.
\]
\end{definition}

Это числовой аналог понятия полуразрешимости.

\subsection{Связь с языками}
Если рассматривать язык как отношение принадлежности слова языку, то:
\begin{itemize}
    \item рекурсивные отношения соответствуют разрешимым языкам;
    \item полурекурсивные отношения соответствуют перечислимым языкам.
\end{itemize}

\section{Алгоритмически неразрешимые задачи}

\subsection{Определение}
\begin{definition}
Задача называется \textbf{алгоритмически неразрешимой}, если не существует алгоритма, решающего её на всех допустимых входах.
\end{definition}

\subsection{Общий факт существования}
Из кардинальных соображений:
\begin{itemize}
    \item алгоритмов счётно много;
    \item задач над натуральными числами или языков несчётно много.
\end{itemize}
Значит, неразрешимые задачи существуют.

\section{Проблема останова}

\subsection{Формулировка}
\begin{definition}
\textbf{Проблема останова}:
\[
K=\{\,x \mid \varphi_x(x)\downarrow \,\}.
\]
То есть по коду программы $x$ нужно определить, завершится ли эта программа на входе $x$.
\end{definition}

\subsection{Почему именно такой вид}
Самоприменение $x$ к $x$ нужно для диагонального рассуждения. Именно оно позволяет построить противоречие.

\subsection{Полурекурсивность проблемы останова}
\begin{theorem}
Множество $K$ полурекурсивно.
\end{theorem}

\begin{proof}
Из теорема Клини о нормальной форме следует:
\[
\varphi_x(x)\downarrow \iff (\exists y)\,T(x,x,y),
\]
где $T$ --- предикат Клини, примитивно рекурсивный, а следовательно рекурсивный.

Значит, условие принадлежности $x$ множеству $K$ выражается через существование такого $y$, что выполняется рекурсивное отношение $T(x,x,y)$. Проекция рекурсивного отношения есть полурекурсивное отношение. Следовательно, $K$ полурекурсивно.
\end{proof}

\subsection{Неразрешимость проблемы останова}
\begin{theorem}
Проблема останова неразрешима.
\end{theorem}

\begin{proof}
Предположим противное: пусть $K$ разрешимо. Тогда и дополнение
\[
\overline K=\{\,x\mid \varphi_x(x)\uparrow\,\}
\]
тоже разрешимо, а значит полурекурсивно.

По нормальной форме полурекурсивных отношений существует такой индекс $i$, что
\[
x\in \overline K \iff (\exists z)\,T(i,x,z).
\]
Но по определению $\overline K$:
\[
x\in \overline K \iff \varphi_x(x)\uparrow.
\]

Подставим $x=i$. Получим
\[
i\in \overline K \iff (\exists z)\,T(i,i,z).
\]
Левая часть означает:
\[
\varphi_i(i)\uparrow.
\]
Правая часть по смыслу предиката Клини означает:
\[
\varphi_i(i)\downarrow.
\]
Следовательно,
\[
\varphi_i(i)\uparrow \iff \varphi_i(i)\downarrow,
\]
что невозможно. Противоречие. Значит, $K$ неразрешимо.
\end{proof}

\subsection{Интерпретация доказательства}
Это классическое диагональное доказательство:
\begin{itemize}
    \item предполагается существование универсального решателя;
    \item затем строится объект, который ведёт себя противоположно предсказанию этого решателя;
    \item отсюда получается противоречие.
\end{itemize}

\subsection{Следствия}
\begin{corollary}
Дополнение $\overline K$ не является полурекурсивным.
\end{corollary}

\begin{proof}
Если бы $\overline K$ было полурекурсивным, то и $K$, и $\overline K$ были бы полурекурсивны. Тогда по ранее доказанной теореме $K$ было бы разрешимо. Но это не так.
\end{proof}

\begin{corollary}
Не существует алгоритма, который по программе и входу всегда определяет, завершится ли вычисление.
\end{corollary}

\section{Теорема Клини о нормальной форме и её роль}

\subsection{Нормальная форма}
В конспекте доказывается важнейший результат:

\begin{theorem}[Клини о нормальной форме]
Для любой частично рекурсивной функции $f$ существует такой индекс $z$ и такая примитивно рекурсивная функция $d$, что
\[
f(x)=d\bigl((\mu y)\,T(z,x,y)\bigr).
\]
\end{theorem}

\subsection{Смысл}
Любая частично рекурсивная функция представляется как:
\begin{enumerate}
    \item поиск кода корректного завершающегося вычисления;
    \item декодирование результата этого вычисления.
\end{enumerate}

Это делает возможным доказательство многих результатов о неразрешимости.

\section{Теорема Райса}

\subsection{Подготовительное замечание}
Теорема Райса относится не к синтаксическим свойствам программ, а к \textbf{семантическим} свойствам, то есть свойствам той функции, которую программа вычисляет.

\subsection{Формулировка}
\begin{theorem}[Райса]
Пусть $\mathcal C\subseteq \mathcal P$ --- некоторое множество частично рекурсивных функций. Рассмотрим множество
\[
A=\{\,x\in\mathbb N \mid \varphi_x \in \mathcal C\,\}.
\]
Тогда множество $A$ рекурсивно тогда и только тогда, когда оно тривиально: либо
\[
A=\varnothing,
\]
либо
\[
A=\mathbb N.
\]
\end{theorem}

\subsection{Идея}
Любое нетривиальное свойство вычисляемой функции алгоритмически неразрешимо.

\subsection{Доказательство}
\begin{proof}
\textbf{Достаточность.}
Если $A=\varnothing$, то его характеристическая функция равна $1$ на всех входах.
Если $A=\mathbb N$, то его характеристическая функция равна $0$ на всех входах.
Обе функции рекурсивны.

\textbf{Необходимость.}
Предположим, что $A$ рекурсивно, но нетривиально. Тогда существуют числа
\[
a\in A,\qquad b\notin A.
\]
Поскольку $A$ рекурсивно, функция
\[
f(x)=
\begin{cases}
b, & x\in A,\\
a, & x\notin A
\end{cases}
\]
является рекурсивной.

Заметим, что для любого $x$:
\[
x\in A \iff f(x)\notin A.
\]
Действительно:
\begin{itemize}
    \item если $x\in A$, то $f(x)=b\notin A$;
    \item если $x\notin A$, то $f(x)=a\in A$.
\end{itemize}

По теореме о рекурсии существует такое число $e$, что
\[
\varphi_{f(e)}=\varphi_e.
\]
Но множество $A$ определяется через \emph{вычисляемую функцию}, а не через конкретный код. Поэтому из равенства функций следует:
\[
e\in A \iff f(e)\in A.
\]
С другой стороны, из свойства функции $f$:
\[
e\in A \iff f(e)\notin A.
\]
Противоречие. Значит, нетривиальное множество индексов такого вида не может быть рекурсивным.
\end{proof}

\subsection{Содержательная интерпретация}
Теорема Райса утверждает, что не существует общего алгоритма, который по программе определяет, например:
\begin{itemize}
    \item вычисляет ли программа тождественно нулевую функцию;
    \item является ли вычисляемая функция всюду определённой;
    \item определена ли функция хотя бы на одном аргументе;
    \item принимает ли функция значение $0$ хотя бы где-нибудь;
    \item совпадает ли вычисляемая функция с фиксированной функцией.
\end{itemize}

Если свойство нетривиально, задача неразрешима.

\subsection{Что не покрывает теорема Райса}
Она не относится к чисто синтаксическим свойствам программ, например:
\begin{itemize}
    \item содержит ли код инструкции определённого типа;
    \item превышает ли длина программы некоторое число;
    \item используется ли в коде конкретный символ.
\end{itemize}

Потому что это свойства \emph{кода}, а не вычисляемой функции.

\section{Проверка рекурсивности как пример применения теоремы Райса}

\subsection{Задача}
Рассмотрим множество
\[
R=\{\,x\mid \varphi_x\in\mathcal R\,\},
\]
то есть множество индексов всюду определённых частично рекурсивных функций.

\subsection{Вывод}
Это свойство:
\begin{itemize}
    \item семантическое;
    \item нетривиальное;
\end{itemize}
следовательно, по теореме Райса оно неразрешимо. Более того, можно доказать, что оно даже не является полурекурсивным.

\subsection{Интерпретация}
Нет алгоритма, который мог бы в общем случае даже только подтверждать, что программа завершается на всех входах.

\section{Примеры алгоритмически неразрешимых проблем}

\subsection{Проблема останова}
Уже рассмотрена подробно:
\begin{itemize}
    \item полурекурсивна;
    \item неразрешима.
\end{itemize}

\subsection{Проблема эквивалентности программ}
Пусть даны две программы. Нужно определить, вычисляют ли они одну и ту же функцию.
Это семантическое нетривиальное свойство пары программ, поэтому такая задача неразрешима.

\subsection{Проверка всюду определённости}
По программе определить, останавливается ли она на всех входах. Это неразрешимо и даже не полуразрешимо.

\subsection{Entscheidungsproblem}
Проблема разрешения для логики первого порядка неразрешима. Это связывается с неразрешимостью проблемы останова.

\subsection{Первая теорема Гёделя о неполноте}
Если бы достаточно сильная формальная система была полной и непротиворечивой, можно было бы разрешать проблему останова. Так как это невозможно, полная и непротиворечивая арифметически выразительная система существовать не может.

\subsection{Проблема соответствий Поста}
Даны два набора слов
\[
A=(a_1,\dots,a_n), \qquad B=(b_1,\dots,b_n).
\]
Требуется определить, существует ли последовательность индексов $i_1,\dots,i_k$, такая что
\[
a_{i_1}a_{i_2}\dots a_{i_k}=b_{i_1}b_{i_2}\dots b_{i_k}.
\]
Общая задача неразрешима.

\subsection{Десятая проблема Гильберта}
Не существует универсального алгоритма решения произвольных диофантовых уравнений.

\section{Сводная сравнительная таблица моделей}

\begin{center}
\begin{longtable}{|>{\raggedright\arraybackslash}p{3.7cm}|>{\raggedright\arraybackslash}p{5.4cm}|>{\raggedright\arraybackslash}p{5.4cm}|}
\hline
\textbf{Модель} & \textbf{Как устроена} & \textbf{Как выполняется шаг} \\
\hline
Нормальный алгоритм Маркова & Упорядоченный набор подстановок над словами & Находим первую применимую подстановку, заменяем первое вхождение, при заключительной подстановке останавливаемся \\
\hline
Машина Тьюринга & Лента, головка, конечное множество состояний, функция переходов & Читаем символ, пишем новый, меняем состояние, двигаемся влево/вправо \\
\hline
Рекурсивные функции & Инициальные функции и операторы построения & Вычисляем функцию через композицию, примитивную рекурсию и, в общем случае, неограниченный поиск \\
\hline
\end{longtable}
\end{center}

\section{Краткие формулы, которые полезно помнить}

\subsection*{Разрешимость и перечислимость}
\[
L \text{ разрешим } \Rightarrow L \text{ перечислим}.
\]
\[
L,\overline L \text{ перечислимы } \Rightarrow L \text{ разрешим}.
\]

\subsection*{Проблема останова}
\[
K=\{x\mid \varphi_x(x)\downarrow\}.
\]
\[
K \in \mathcal P^*, \qquad K \notin \mathcal R^*.
\]

\subsection*{Теорема Райса}
Для нетривиального $\mathcal C\subseteq\mathcal P$:
\[
A=\{x\mid \varphi_x\in \mathcal C\}
\]
не является рекурсивным.

\subsection*{Частично рекурсивные функции}
\[
\mathcal P = \text{замыкание инициальных функций относительно композиции, примитивной рекурсии и }\mu\text{-оператора}.
\]

\section{Заключение}
\begin{enumerate}
    \item Понятие алгоритма изначально интуитивно, но для строгой математики оно должно быть формализовано.
    \item Такими формализациями служат, в частности, нормальные алгоритмы Маркова, машины Тьюринга и рекурсивные функции.
    \item Все эти модели эквивалентны: они задают один и тот же класс вычислимых объектов.
    \item Это позволяет определить разрешимые и перечислимые языки, а также частично и всюду вычислимые функции.
    \item Не все задачи алгоритмически разрешимы.
    \item Базовый пример неразрешимой задачи --- проблема останова.
    \item Ещё более общий результат --- теорема Райса: любое нетривиальное свойство вычисляемой функции неразрешимо.
\end{enumerate}

\end{document}
